\section{Steady-state flux reconstruction}
Begin with the two voltage balances
\begin{align}
 \ud&=\Rs\id-\omegae\psiq, &
 \uq&=\Rs\iq+\omegae\psid.\label{eq:voltage}
\end{align}
They apply when the rotor-frame flux derivative vanishes, the represented
voltage is the terminal fundamental voltage, and the lumped resistive model
is adequate. Set
\begin{equation}
 \voltage=\begin{bmatrix}\ud\\\uq\end{bmatrix},\quad
 \current=\begin{bmatrix}\id\\\iq\end{bmatrix},\quad
 \flux=\begin{bmatrix}\psid\\\psiq\end{bmatrix},\quad
 \rotationgenerator=\begin{bmatrix}0&-1\\1&0\end{bmatrix}.
\end{equation}
Then $\rotationgenerator\flux=[-\psiq,\psid]\trans$ rotates a vector counterclockwise
through $90^\circ$, and \cref{eq:voltage} becomes
\begin{equation}
 \voltage=\Rs\current+\omegae\rotationgenerator\flux.\label{eq:vector-voltage}
\end{equation}
Subtracting the resistive drop isolates rotational EMF. Rotating it back
undoes its quadrature orientation. Since $\rotationgenerator^2=-\matr I$, for
$\omegae\ne0$ this gives
\begin{equation}
 \boxed{\flux^{\rm rec}=-\frac{1}{\omegae}\rotationgenerator
 (\voltage-\Rs\current).}\label{eq:reconstruction}
\end{equation}
Equivalently,
\begin{equation}
 \psid^{\rm rec}=\frac{\uq-\Rs\iq}{\omegae},\qquad
 \psiq^{\rm rec}=\frac{\Rs\id-\ud}{\omegae}.
\end{equation}
This operation reveals flux from the voltage balance but hides the origin of
the subtracted voltage. Any error in voltage or resistive drop is divided by
speed. At zero speed the stationary equation contains no flux information;
near zero speed the inversion amplifies noise.

No reflection symmetry follows from this inversion. Indeed, any chosen
flux-current function, symmetric or otherwise, can be inserted in
\cref{eq:vector-voltage} to produce corresponding voltages. Structural
symmetry must come from the machine, not from solving its voltage equation.
The alternative componentwise derivation gives the same signs: $u_q-R_si_q$
contains $+\omega_e\psi_d$, whereas $u_d-R_si_d$ contains
$-\omega_e\psi_q$.
